\documentclass[12pt]{article}
\usepackage{booktabs}
\usepackage{caption}
\usepackage{threeparttable}
\usepackage{dcolumn}
\usepackage{adjustbox}
\usepackage[utf8]{inputenc}
\usepackage[margin=0.5in, paperwidth=14in, paperheight=10in]{geometry}
\pagestyle{empty}
\begin{document}
\setcounter{table}{6}
\begin{table}[htbp]
\centering
\begin{threeparttable}

\caption{Ranking de departamentos}
\begin{tabular}[t]{lrrlrr}
\toprule
Departamento & Costo Mol & IVIG & Grupo & Productividad & PIB\\
\midrule
CUNDINAMARCA & 1219.241 & 0.745 & Vulnerable & 153.282 & 99.945\\
TOLIMA & 1162.650 & 0.735 & Vulnerable & 141.333 & 33.735\\
CAUCA & 1158.889 & 0.702 & Vulnerable & 150.768 & 28.537\\
VALLE DEL CAUCA & 1163.569 & 0.679 & Vulnerable & 152.836 & 153.564\\
BOGOTA, D.C. & 1080.999 & 0.522 & Vulnerable & 146.220 & 395.438\\
\addlinespace
CALDAS & 1234.233 & 0.511 & Control & 190.927 & 26.213\\
ANTIOQUIA & 997.440 & 0.501 & Control & 154.082 & 231.122\\
RISARALDA & 1312.114 & 0.450 & Vulnerable & 145.678 & 26.404\\
BOLIVAR & 1322.401 & 0.316 & Control & 162.806 & 56.924\\
SANTANDER & 1028.903 & 0.231 & Control & 177.290 & 101.446\\
\addlinespace
BOYACA & 1098.967 & 0.182 & Control & 179.730 & 42.277\\
ATLANTICO & 1458.762 & 0.176 & Vulnerable & 173.625 & 70.862\\
CORDOBA & 1021.598 & 0.050 & Control & 144.475 & 28.303\\
\bottomrule
\end{tabular}
\begin{tablenotes}
\footnotesize
\item \textit{Notas:}El PIB departamental en Billones COP
\end{tablenotes}
\end{threeparttable}
\end{table}
\end{document}
